\chapter{Introduction}
\label{c:introduction}

% Struktur
\section{Motivation}
Mobile applications (apps) can have a variety of vulnerabilities. Modern Android apps have, on average, a surprisingly large number of vulnerabilities, which greatly impacts their security and usage \todo{Add any of the papers read here that say this}. Because of this, it is of great importance to have tools that can detect vulnerabilities, both statically by inspecting the code and dynamically by executing the code and inspecting it at runtime. There exists a diverse landscape of such Static Security Application Testing (SAST) and Dynamic Security Application Testing (DAST) tools. But how can the effectiveness and accuracy of these tools be measured and compared to one another? 

To measure the effectiveness of a tool, it is just as important to know how many security vulnerabilities the tool was not able to detect, as it is important to find the vulnerabilities themselves. Since most real-world mobile apps are closed source \todo{Probably add source for this}, there is no ground truth against which the performance of a security analysis tool on such apps can be evaluated, making it impossible to use them for measuring the effectiveness of such tools. To effectively measure and compare the accuracy of existing analysis tools, they have to be evaluated on intentionally vulnerable apps, of which the existing vulnerabilities are known to the person conducting the tests. 

\section{Problem statement}
The intentionally vulnerable apps dataset hosted by OWASP as a reference for its Mobile Application Testing Guide (MASTG) \cite{owasp_reference_app_website} is a very valuable learning resource, teaching mobile developers how security vulnerabilities can look like and how to find them. These intentionally vulnerable apps are furthermore good resources for benchmarking mobile security tools. 

Many of these MASTG reference apps have unfortunately not been updated in years and are designed for older versions of Android, meaning that they do not adequately reflect and implement the modern vulnerabilities discussed by OWASP in its testing guide or weakness enumeration anymore. This greatly diminishes their educational value as teaching resources. Most of them do not contain documentation of their implemented vulnerabilities, which negatively impacts their usefulness for benchmarking security tools, as it makes it more difficult to assess how many vulnerabilities have not been found by the tools. Because of this, there is a gap between the educational value provided by the OWASP MAS project and how it is reflected in its MASTG reference apps and the practical value it provides in mobile security testing and research. 

\section{Objectives of the Thesis}
The objectives of this thesis largely concern developing multiple open-source OWASP MASTG reference apps that include both static and dynamic security vulnerabilities. To achieve this, the thesis pursues the following objectives:

% Aufzählung
\begin{enumerate}
    \item Analyze and understand the OWASP MAS ecosystem and its vulnerability categories.
    \item Conduct a literature review to better understand how OWASP MASTG, its reference apps, as well as other related components, are used in the academic world.
    \item Evaluate the state-of-the-art in intentionally vulnerable OWASP MASTG Android reference apps with respect to coverage, documentation, educational value and usability for security tool benchmarking.
    \item Design and implement new MASTG Android reference apps, adhering to the determined state-of-the-art standards. The focus lies on vulnerabilities which have seen little to no coverage in existing reference apps. Both static and dynamic vulnerabilities are to be included. 
    \item Document each implemented vulnerability by mapping it to its MASWE counterpart, explaining where it can be found in the implementations code, and how it can be exploited.
    \item Compare our reference apps with the existing reference apps landscape and assess the contribution made by our apps. Compare the coverage of vulnerabilities with and without our apps. Evaluate our apps for suitability as teaching resource and for benchmarking security analysis tools.
\end{enumerate}

\section{Methodology and Structure of the Thesis}

A design and implementation focused research approach is used for this thesis. It consists of first reviewing and analyzing the existing work in the related field, and is then followed by the design, implementation, and evaluation of multiple pieces of software. The thesis begins with an analysis of the OWASP MASTG and its reference applications. Based on the results of this analysis, the choices of which vulnerabilities are implemented are made and multiple new reference apps are designed and developed. The resulting applications are then compared to the existing catalog of MASTG reference apps, and evaluated for their coverage of vulnerabilities, as well as for their suitability for educational purposes and the benchmarking of security analysis tools.

This thesis is structured as follows. Chapter~1 introduces the subject of the thesis, the motivation behind it, its objectives, and the method for achieving them. In Chapter~2, we discuss the methodology \todo{Maybe remove mention of methodology here} and findings of both the literature review and the analysis of the state-of-the-art in existing OWASP MASTG reference apps. Chapter~3 provides essential background knowledge on key Android concepts and the OWASP MAS ecosystem, which are discussed and used in this thesis. Chapter~4 discusses the design for developing the reference apps, including design decisions regarding vulnerability selection as well as system architecture. In Chapter~5, each vulnerability implemented is covered in greater detail, supplemented with code snippets, implementation decisions, and challenges faced during their development. Chapter~6 concerns itself with the results of our work, comparing our developed reference apps with the existing MASTG catalog and assessing the contributions made. Chapter~7 discusses possible threats to the validity of our work, as well as any problems left open or newly discovered. In Chapter~8, we look forward to future work, outlining how this thesis could be extended upon. Chapter~9 then closes out the thesis.